\chapter{先学会安全地走一步}\label{sec:descent}
\section{平方和下降，到底还缺什么}
令 $E(x)=\sum_i x_i^2$。平均一组 $S$，其均值为 $a$，则
\[
E(x)-E(T_Sx)=\sum_{i\in S}x_i^2-pa^2
=\sum_{i\in S}(x_i-a)^2.
\]
右边非负，且除非选中的值全相等，否则严格为正。这是“平均让数更接近”的精确版本。

要从这里得到有限终止，还缺两个条件。第一，中间数最好都是整数，这样下降量不可能无穷趋近于零。第二，别走到上一讲的共同余数陷阱里。我们将把这两个条件和一种位置储备一起保持：\textbf{有两个不同数值，各出现至少 $p$ 次。}

这种数值叫“重值”。比如五平均中 $a$ 出现六次，它是重值；出现四次则不是。重值不是说 $a$ 的绝对值大。

\section{留在外面的两个数，为什么能保护这一步}
假如未选位置里还留着一个奇数和一个偶数，那么不管选中的位置怎样变化，输出都不可能全同奇偶。模任意素数也是同一个道理。因此，与其直接检查新的均值，不如事先在外面保留一对不同余位置。

什么时候这会紧张？只有绝大多数位置都属于同一个余数类、少数例外很容易被一次全拿走的时候。
\begin{definition}[危险素数与危险集]
称 $\ell\mid n$ 且 $\ell\ne p$ 为危险素数。
若模 $\ell$ 有一个余数类包含至少 $n-p$ 个位置，称它为主类，
将其余位置构成的集合记为 $F_\ell$。
如果一次平均会使全部位置同余，那么未选的 $n-p$ 个位置已经必须同余。
因此只有这种主类存在时，才有一步制造障碍的风险。
\end{definition}
例如五平均二十元中，若十八个位置是偶数、两处是奇数，那么不要把两处奇数同时拿走；保留一处，就保留了奇偶差异。对每个危险素数留一个例外位置，总共最多留 $d_p(n)$ 个。
\begin{lemma}[由未选位置保证安全]\label{lem:safety}
设 $n>2p$。若模危险素数 $\ell$ 存在大小至少 $n-p$ 的余类，其补集 $F_\ell$ 唯一，且 $2\le|F_\ell|\le p$。不包含任何完整 $F_\ell$ 的整数平均安全。以下整数选择也安全：
\begin{enumerate}
\item $(a^{p-1},x)$，其中至少一份 $a$ 未选；
\item 由若干份 $a,b$ 及一份第三值 $x$ 组成，且 $a,b$ 均有未选副本；
\item 每种被选数值均有未选副本。
\end{enumerate}
\end{lemma}
\begin{proof}
$n-p>n/2$ 给出主类唯一性。例外数不能为 $1$，否则零和模 $\ell$ 迫使唯一例外仍属主类。避开完整 $F_\ell$ 时，未选位置保留主类及例外。对后三种选择，若输出全同余，未选副本及选中和为 $p$ 倍新均值依次迫使全部被选值也属于共同余类，矛盾。
\end{proof}
这里还给了几个更省位置的选择。最常用的是 $(a^{p-1},x)$：留下的一份 $a$ 会迫使任何假想的共同余数等于 $a$，而平均方程再迫使 $x$ 也等于这个余数，矛盾。正是这一份重复值，省掉了另找见证的麻烦。

\section{为什么要保留两个重值}
一次平均会产生 $p$ 份相等的新值。如果另外还保留着一个旧的 $p$ 重值，而且新旧两值不同，就又回到“两个重值”的状态，可以继续下一轮。

注意，新均值若恰好等于保留的旧值，就只剩一个重值。这种情况叫“碰撞”。它不是说原配置不可解，只是当前的储备没有保持住。下一讲会通过碰撞等式重新选一组。

初始化时，在分析中先考虑 $pX$。所有初值都被 $p$ 整除，所以第一次任取 $p$ 项的均值是整数。选出位置序列后，在原 $X$ 上执行同一序列即可。后面的整数尺度固定下来，不在下降途中随意改变。

一个小技巧反复有用：对 $p+1$ 个候选数，删掉一个后剩 $p$ 个。如果所有删法都得到同一个和，比较 $S-x_i$ 与 $S-x_j$ 就知 $x_i=x_j$。因此，多留一个候选位置便有机会避开指定均值。
\begin{lemma}[建立两个不同重值]\label{lem:initialize}
若 $n\ge3p+d$，从 $pX$ 出发至多两次整数平均可建立两个不同重值。
\end{lemma}
\begin{proof}
在任一候选池中，为仍完整包含在池内的每个 $F_\ell$ 留出一个位置，至多留出 $d$ 个。若池内所有值同余模 $p$、每种实际值不足 $p$ 份，余下 $p$ 个位置给安全非恒等整数平均；余下 $p+1$ 个位置时还可避开一个指定均值，因为所有删一子集同和会迫使这 $p+1$ 个值全相等。
已有两个重值时无须操作。仅一个重值 $a$ 且重数至少 $2p-1$ 时，选择 $(a^{p-1},b)$，留下 $p$ 份 $a$ 并生成另一重值。若重数至多 $2p-2$，在 $a$ 外的池中避开均值 $a$。若初始没有重值，先制造一个；它的重数至多 $2p-1$，其外至少 $d+p+1$ 个未改变的初值仍被 $p$ 整除，故第二次可制造不同重值。
\end{proof}

\section{整除要求交给一个零和定理}
安全性安排好以后，还要使所选和被 $p$ 整除。只看模 $p$ 的余数，这就是找 $p$ 项零和。Erdős--Ginzburg--Ziv 定理在此处的结论很直观：$2p-1$ 个候选永远够；少一个时，失败的形状也能描述。
\begin{example}[五平均的边界池]
模 $5$ 的池 $(0^4,1^4)$ 没有五项零和。
五项中必须包含 $1,2,3$ 或 $4$ 个 $1$，和都不被 $5$ 整除。
添加任意一个第九位置以后，零和引理保证出现五项零和。
逆命题说，长度为八的失败池在仿射改名后只能有这种两类各四份的形状。
\end{example}
\begin{lemma}[素数零和选择及逆边界]\label{lem:egz}
任意 $2p-1$ 个 $\F_p$ 元素包含 $p$ 项零和。长度 $2p-2$ 的序列没有 $p$ 项零和，当且仅当它为 $u^{p-1}v^{p-1}$、$u\ne v$。
\end{lemma}
\begin{proof}
先说明有限域计数：若 $m$ 元多项式 $F_i$ 的次数和小于 $m$，其共同零点数模 $p$ 等于
\[\sum_{z\in\F_p^m}\prod_i(1-F_i(z)^{p-1})=0.
\]
这里把 $1-F_i(z)^{p-1}$ 看成一个指示函数：$F_i(z)=0$ 时它为 $1$，否则为 $0$。
对一个变量，有 $\sum_{t\in\F_p}t^k=0$，除非 $k$ 是正的 $p-1$ 倍数；当 $k=0$ 时求和也是 $p=0$。
展开后每个单项式的总次数小于 $m(p-1)$，至少一个变量的指数不是正的 $p-1$ 倍数，对它求和为零。应用于
\[\sum_{i=1}^{2p-1}z_i^{p-1}=0,\qquad
\sum_{i=1}^{2p-1}a_i z_i^{p-1}=0.
\]
由于 $z_i^{p-1}$ 在 $z_i\ne0$ 时为 $1$，第一个方程要求非零坐标个数被 $p$ 整除。
一共只有 $2p-1$ 个变量，所以非零解的支持恰有 $p$ 项；第二个方程要求该支持对应的 $a_i$ 之和为零。
每个合格支持允许每一非零坐标任取 $p-1$ 个值，零点数为 $1+(p-1)^p M$，故 $M\equiv1\pmod p$。

证明逆命题。出现 $p$ 次同一余数时已有零和；仅两类时只能为所示形式。若至少三类且各重数至多 $p-1$，可选两对不交位置，使两对的差非零且不互为正负。取重数最大的类 $a$，它至少有两份；若两个其他类对 $a$ 的差不互为正负即可。否则仅有 $a,b,c$ 三类，且 $b,c$ 关于 $a$ 对称；$b,c$ 至少一类有两份，改用 $(a,b),(b,c)$，两差相差倍率 $2$。
余下 $2p-6$ 项最大重数至多 $p-2$，可再配 $p-4$ 对不同余数，留下两个单点。所用匹配事实为：多重数序列的异类最大匹配数是 $\min(\lfloor M/2\rfloor,M-f_{\max})$。若某类超过一半，将其与其余类配对；否则将各类连续排成一列，前后半列配对即达到该界。
前两对各取一项的和集有四个元素。以后每加一对，和集至少增大一，除非已为 $\F_p$：一个真子集不能在非零平移下不变。故 $p-2$ 对的选择和遍历 $\F_p$，加上两个单点得到 $p$ 项零和。
\end{proof}

第一次读这个证明，可以把它拆成两幅图景。前半段用 $z_i^{p-1}$ 表示“第 $i$ 项选不选”：它在零处为零，其余地方为一。第一方程迫使选项数是 $p$，第二方程迫使所选余数之和是零。有限域求和证明这样的选择确实存在。

后半段处理少一个候选的边界。把一些位置两两配对，每对可以选左边或右边。这提供一串二选一的自由度。每增加一对，可能的余数和至少增加一种；只要一开始比普通情况多出一种，就能覆盖全部 $p$ 个余数。唯一躲开这个过程的，是两类各 $p-1$ 份的形状。

到此，我们有了三件工具：留下余数见证保证安全；重复值保证下一轮还有材料；零和定理保证平均仍为整数。下一讲只负责把它们安排成一个覆盖所有状态的选择方法。
